\documentclass[journal]{IEEEtran}
\IEEEoverridecommandlockouts
\usepackage{cite}
\usepackage{amsmath,amssymb,amsfonts}
\usepackage{algorithmic}
\usepackage{graphicx}
\usepackage{textcomp}
\usepackage{xcolor}
\usepackage{hyperref}
\usepackage{booktabs}
\usepackage{multirow}

\def\BibTeX{{\rm B\kern-.05em{\sc i\kern-.025em b}\kern-.08em
    T\kern-.1667em\lower.7ex\hbox{E}\kern-.125emX}}

\title{Ollama-Optimizer: Automated Quantization-Aware Optimization for Large Language Model Deployment on Consumer Hardware}

\author{\IEEEauthorblockN{A. Mahesh Bhat}
\IEEEauthorblockA{\textit{Independent Researcher}\\
Bangalore, India\\
M.Tech Advanced Computing, NIT Bhopal, 2023\\
Email: mbhat648@gmail.com}}

\begin{document}

\maketitle

\begin{abstract}
Large Language Models (LLMs) have revolutionized natural language processing but pose significant deployment challenges due to their enormous memory requirements and computational demands. A 70-billion parameter model at FP16 precision requires approximately 140 GB of RAM, rendering it impractical for deployment on consumer hardware. Post-training quantization (PTQ) addresses this challenge by compressing model weights from 16-bit floating-point to lower-precision integer representations (INT8, INT4, or lower), reducing memory footprint by up to 75\% while maintaining acceptable quality. However, selecting optimal quantization parameters requires expertise in both hardware architecture and quantization theory, creating a barrier for widespread adoption. We present \textbf{Ollama-Optimizer}, an automated system for hardware-aware LLM optimization that automatically detects system capabilities, analyzes installed models, recommends optimal quantization configurations, and benchmarks performance improvements. Our system employs a multi-stage optimization pipeline: (1) hardware profiling across CPU, RAM, and GPU/VRAM on Windows, Linux, and macOS; (2) model analysis parsing GGUF quantization metadata; (3) optimization planning using quantization-aware memory calculations; (4) model creation via Ollama Modfile generation; and (5) performance benchmarking measuring tokens/sec, latency, and memory consumption. Evaluation on an Apple M2 system (4.0 GB VRAM, 1.2 GB available RAM) with six installed models (gemma3:27b, glm-5:cloud, qwen3:8b, nomic-embed-text:latest, llama2:latest, llama3:70b) demonstrates successful hardware detection and quantization recommendations. Benchmarking of llama2:latest (7B parameters, Q4\_0 quantization, 3.6 GB) achieved 18.3 tokens/sec average throughput with 302 ms time-to-first-token (TTFT) and 6 MB memory consumption. The system correctly identified that large models (gemma3:27b at 16.2 GB and llama3:70b at 37.2 GB) exceed available memory capacity and recommended alternative quantization strategies or smaller models. Ollama-Optimizer provides a unified, cross-platform solution for democratizing LLM deployment by automating the complex optimization decisions traditionally requiring expert knowledge.
\end{abstract}

\begin{IEEEkeywords}
large language models, quantization, post-training quantization, GGUF, hardware-aware optimization, consumer hardware deployment, model compression, inference acceleration
\end{IEEEkeywords}

\section{Introduction}

The emergence of Large Language Models (LLMs) such as GPT-3 \cite{brown2020}, LLaMA \cite{touvron2023}, and Mistral \cite{jiang2023} has transformed natural language processing, enabling unprecedented capabilities in text generation, reasoning, and task completion. However, these models pose significant deployment challenges due to their enormous size: a 70-billion parameter model at FP16 precision requires approximately 140 GB of RAM \cite{frantar2022}, far exceeding the capabilities of consumer hardware. This memory barrier restricts LLM deployment to well-funded organizations with access to high-performance computing infrastructure, limiting widespread adoption and on-premise deployment.

Quantization addresses this challenge by compressing model weights from high-precision floating-point (FP32, FP16) to lower-precision integer representations (INT8, INT4, or lower). Post-training quantization (PTQ) techniques such as GPTQ \cite{frantar2022}, AWQ \cite{lin2023}, and SmoothQuant \cite{xiao2022} can reduce model size by up to 75\% while maintaining acceptable quality. The GGUF format (successor to GGML) developed by Georgi Gerganov for the llama.cpp project \cite{gerganov2023} has emerged as a de facto standard for quantized LLM distribution, supporting multiple quantization schemes (Q8\_0, Q4\_K\_M, Q2\_K, etc.) with varying trade-offs between memory footprint, inference speed, and quality.

Despite these advances, selecting optimal quantization parameters remains challenging. The decision involves: (1) understanding hardware constraints (CPU cores, RAM, GPU/VRAM), (2) calculating memory requirements at different quantization levels, (3) balancing quality retention against speed improvements, (4) determining optimal GPU layer offloading, and (5) tuning runtime parameters (threads, context window, batch size). These decisions require expertise in both hardware architecture and quantization theory, creating a barrier for non-expert users.

We present \textbf{Ollama-Optimizer}, an automated system for hardware-aware LLM optimization that democratizes access to quantization techniques. Our system automatically detects system capabilities, analyzes installed models, recommends optimal quantization configurations, generates optimized model variants, and benchmarks performance improvements. Built on concepts from ngrok's ``Quantization from the Ground Up'' blog post \cite{ngrok2023}, Ollama-Optimizer implements a comprehensive quantization knowledge base and optimization pipeline that requires no expert knowledge from the end user.

\textbf{Contributions:}
\begin{itemize}
    \item Cross-platform hardware profiling for CPU, RAM, GPU/VRAM detection across Windows, Linux, and macOS with support for NVIDIA, AMD, and Apple Silicon GPUs
    \item Automated model analysis parsing GGUF quantization metadata to determine current quantization level, parameter count, and family
    \item Quantization-aware optimization planning that calculates memory requirements at all quantization levels and recommends configurations fitting available hardware
    \item Priority-based optimization strategies (quality, speed, balanced, minimum) allowing users to trade off quality against memory and speed
    \item Integrated benchmarking measuring tokens/sec, time-to-first-token (TTFT), prompt processing rate, and memory consumption before and after optimization
    \item Built-in quantization education with interactive explainer demonstrating absmax, zero-point, and per-group quantization techniques
\end{itemize}

\section{Related Work}

\subsection{Existing Quantization Tools}

Several tools exist for LLM quantization and optimization, but each has limitations that Ollama-Optimizer addresses:

\textbf{llama.cpp} \cite{gerganov2023} provides the core quantization engine with GGUF format support but requires manual command-line invocation for quantization (`llama-quantize`) and does not provide hardware-aware recommendations. Users must manually determine optimal quantization levels based on their hardware, requiring expertise in both quantization theory and system architecture.

\textbf{AutoGPTQ} \cite{frantar2022} is a Python library implementing the GPTQ algorithm but focuses only on the quantization algorithm itself, not on deployment optimization. It does not provide hardware detection, GPU layer offloading recommendations, or runtime parameter tuning.

\textbf{GPTQ-for-LLaMA} extends GPTQ specifically for LLaMA models but requires manual configuration of quantization parameters and does not integrate with deployment platforms like Ollama.

\textbf{Ollama} provides model management and deployment but does not include optimization capabilities. Users must manually determine which quantization level to use and how to configure runtime parameters for their hardware.

Ollama-Optimizer differs from these tools by providing an end-to-end automated pipeline that integrates hardware detection, quantization planning, model creation, and benchmarking in a single system. Unlike manual tools, Ollama-Optimizer requires no user expertise in quantization theory or hardware architecture.

\subsection{Quantization Techniques for LLMs}

\textbf{Post-Training Quantization (PTQ)}: GPTQ \cite{frantar2022} proposes accurate post-training quantization for generative pre-trained transformers using second-order information to minimize quantization error. GPTQ quantizes 1--3 billion parameter models in minutes and 175B models in hours. AWQ (Activation-aware Weight Quantization) \cite{lin2023} identifies that only 1\% of weights are critical for accuracy and proposes activation-aware weight quantization preserving these salient weights. SmoothQuant \cite{xiao2022} proposes equivalent transformations that smooth activation magnitudes for easier quantization. These methods achieve 3--4 bit quantization with minimal accuracy loss but require expertise to apply correctly.

\textbf{Quantization-Aware Training (QAT)}: Quantization-aware training \cite{jacob2018} simulates quantization during training to reduce accuracy loss. While effective, QAT requires access to training data and compute resources, making it impractical for most users. PTQ methods are preferred for deployment scenarios where only the trained model is available.

\textbf{GGUF and llama.cpp}: The llama.cpp project \cite{gerganov2023} pioneered LLM inference on consumer hardware using pure C/C++ implementation. GGUF (Georgi's Unified Format) bundles model weights, tokenizer data, and metadata in a single file, supporting multiple quantization schemes. Q4\_K\_M (4-bit K-quant mixed precision) has emerged as a practical sweet spot, offering ~85\% quality at 30\% of FP16 memory \cite{llamacpp2023}.

\subsection{Hardware-Aware Optimization}

\textbf{GPU Layer Offloading}: llama.cpp supports GPU acceleration through layer offloading, where attention layers are computed on GPU while remaining layers execute on CPU \cite{llamacpp2023}. The optimal number of GPU layers depends on VRAM capacity and model size. Ollama-Optimizer automatically calculates optimal offloading based on detected VRAM.

\textbf{Runtime Parameter Tuning}: Inference performance depends on thread count, context window size, and batch size. Ollama-Optimizer automatically tunes these parameters based on CPU cores, available RAM, and model characteristics.

\subsection{Quantization Evaluation}

\textbf{Quality Metrics}: Perplexity is commonly used to measure quantization quality \cite{frantar2022}, but recent work shows perplexity may be misleading for instruction-tuned models \cite{liu2024}. Alternative metrics include task-specific accuracy benchmarks (MMLU, HellaSwag) and KL divergence \cite{llamacpp2023}. Ollama-Optimizer reports estimated quality retention based on quantization level while acknowledging limitations of perplexity-based evaluation.

\textbf{Performance Metrics}: Inference performance is measured by tokens/sec (throughput), time-to-first-token (latency), prompt processing rate, and memory consumption \cite{lin2023}. Ollama-Optimizer benchmarks these metrics using diverse prompts (QA, reasoning, creative, code, summarization).

\subsection{System Architecture}

\subsection{Overview}

Ollama-Optimizer implements a five-stage pipeline: (1) System Profiling, (2) Model Analysis, (3) Optimization Planning, (4) Model Creation, and (5) Benchmarking. The system is implemented in Python 3.9+ using Click for CLI, Rich for terminal reports, psutil for hardware detection, and requests for Ollama REST API communication. The codebase consists of seven modules: cli.py (command-line interface), system\_profiler.py (hardware detection), ollama\_client.py (Ollama API wrapper), quantization.py (quantization knowledge base), optimizer.py (optimization logic), benchmark.py (performance measurement), and reporter.py (rich terminal output).

\subsection{Hardware Profiling}

The SystemProfiler module detects hardware capabilities across platforms:
\begin{itemize}
    \item \textbf{CPU}: Core count, thread count, architecture (x86, ARM), and frequency using psutil.cpu\_count(logical=False) for physical cores and psutil.cpu\_count(logical=True) for logical threads
    \item \textbf{RAM}: Total memory, available memory using psutil.virtual\_memory().total and psutil.virtual\_memory().available
    \item \textbf{GPU}: NVIDIA GPU detection via subprocess.run(['nvidia-smi', '--query-gpu=name,memory.total', '--format=csv,noheader']), AMD GPU via rocm-smi, Apple Silicon GPU via Metal framework using pyobjc
    \item \textbf{VRAM}: GPU memory capacity and available VRAM parsed from nvidia-smi output or Metal API
    \item \textbf{Disk}: Available disk space for model storage using shutil.disk\_usage()
\end{itemize}

Cross-platform support includes Windows (nvidia-smi), Linux (nvidia-smi, rocm-smi), and macOS (Metal). Flash Attention support is detected for Apple Silicon by checking for the presence of the flash-attn package via importlib.util.find\_spec('flash\_attn').

\subsection{Model Analysis}

The OllamaClient module queries the Ollama REST API (http://localhost:11434) to:
\begin{itemize}
    \item List all installed models via GET /api/tags
    \item Retrieve model metadata including size, quantization, parameter count, and family via GET /api/show
    \item Parse GGUF quantization levels from model tags (Q8\_0, Q4\_K\_M, Q2\_K, etc.) using regular expression pattern matching
\end{itemize}

The Quantization module maintains a knowledge base of all GGUF quantization levels with:
\begin{itemize}
    \item Bits per weight (e.g., Q4\_K\_M = 4.83 bits)
    \item Memory multiplier (relative to FP16, e.g., Q4\_K\_M = 0.30x)
    \item Estimated quality retention (e.g., Q4\_K\_M = 85\%)
    \item Speed multiplier (e.g., Q4\_K\_M = 2.3x)
\end{itemize}

\subsection{Optimization Planning}

The Optimizer module calculates optimal configurations for each model using the following algorithm:

\begin{algorithmic}
\STATE Input: Model parameters $P$, available RAM $R_{avail}$, available VRAM $V_{avail}$, priority strategy
\STATE For each quantization level $q$ in knowledge base:
\STATE Calculate memory requirement: $M_q = P \times Bits_q \times 1.2$
\STATE If $M_q \leq (R_{avail} + V_{avail})$: Add $q$ to feasible set $F$
\STATE Sort $F$ by priority (quality: highest bits, speed: lowest bits, balanced: Q4\_K\_M, minimum: lowest bits fitting)
\STATE Select optimal quantization $q_{opt} = F[0]$
\STATE Calculate GPU layers: $L_{gpu} = \min(L_{total}, \lfloor V_{avail} / S_{layer} \rfloor)$ where $S_{layer}$ is approximate layer size
\STATE Set threads: $T = \min(C_{cores}, 16)$
\STATE Set context: $C = \min(4096, \lfloor R_{avail} / 1000 \rfloor)$
\STATE Output: $(q_{opt}, L_{gpu}, T, C)$
\end{algorithmic}

Priority strategies:
\begin{itemize}
    \item \textbf{Quality}: Select highest quantization level fitting in memory (prioritizes accuracy)
    \item \textbf{Speed}: Select most aggressive quantization (lowest bits, maximizes throughput)
    \item \textbf{Balanced}: Select Q4\_K\_M (85\% quality, 2.3x speed) if available, otherwise nearest equivalent
    \item \textbf{Minimum}: Select quantization fitting in minimal hardware (prioritizes memory efficiency)
\end{itemize}

\subsection{Model Creation}

The system generates an Ollama Modfile with optimized configuration:
\begin{verbatim}
FROM <base_model>
PARAMETER num_gpu <gpu_layers>
PARAMETER num_thread <threads>
PARAMETER num_ctx <context>
\end{verbatim}

The optimized model is created via `ollama create <new_tag> -f <modfile>` and tagged with `-optimized` suffix. The Modfile template is generated dynamically based on the calculated parameters.

\subsection{Benchmarking}

The Benchmark module runs 5 diverse prompts covering different task types:
\begin{itemize}
    \item \textbf{QA}: ``What is the capital of France?'' (factual knowledge retrieval)
    \item \textbf{Reasoning}: ``If all cats are animals and all animals are mortal, are all cats mortal?'' (logical reasoning)
    \item \textbf{Creative}: ``Write a short poem about technology'' (creative generation)
    \item \textbf{Code}: ``Write a function to reverse a string in Python'' (code generation)
    \item \textbf{Summarization}: ``Summarize the following text: [long text]'' (text summarization)
\end{itemize}

Metrics measured:
\begin{itemize}
    \item \textbf{Tokens/sec}: Generation throughput calculated as $N_{tokens} / T_{generation}$
    \item \textbf{Time-to-first-token (TTFT)}: Latency before first token, measured from request start to first token received
    \item \textbf{Prompt tokens/sec}: Input processing rate calculated as $N_{prompt} / T_{processing}$
    \item \textbf{Memory}: RAM consumption measured via psutil.Process().memory\_info().rss
\end{itemize}

Each prompt is run twice to account for variability. Results are displayed in Rich-formatted tables with before/after comparison when optimization is successful.

\section{Experimental Evaluation}

\subsection{Methodology}

Our experimental evaluation follows a systematic methodology to assess the effectiveness of Ollama-Optimizer across hardware detection, model analysis, optimization planning, and benchmarking capabilities. The evaluation was conducted on a single test system with the following procedure:

\textbf{Phase 1: System Setup.} We installed Ollama (latest version) and Ollama-Optimizer v1.0.0 on an Apple M2 system running macOS. The system was configured with default Ollama settings, and six models were pre-installed representing various parameter scales and quantization levels.

\textbf{Phase 2: Hardware Detection Validation.} We executed `ollama-optimizer scan` to verify hardware detection accuracy. Detected specifications (CPU cores, RAM, GPU, VRAM) were cross-validated against system information obtained via macOS System Information and Activity Monitor applications.

\textbf{Phase 3: Model Analysis Verification.} We compared model analysis output against ground truth obtained via `ollama list` command, verifying parameter counts, quantization levels, and file sizes. Parsing accuracy was calculated as the ratio of correctly parsed models to total installed models.

\textbf{Phase 4: Memory Constraint Testing.} We tested the system's ability to identify models exceeding available memory by analyzing the recommendations for models larger than 5.2 GB (the sum of available RAM and VRAM). Correct identification was verified by manual calculation of memory requirements using the formula: $Memory = Params \times BitsPerWeight \times 1.2$.

\textbf{Phase 5: Benchmark Execution.} For models fitting in available memory, we executed `ollama-optimizer benchmark --model <name>` with default settings (5 prompts, 2 runs per prompt). We recorded tokens/sec, time-to-first-token (TTFT), prompt processing rate, and memory consumption. Benchmarking was repeated for llama2:latest (3.6 GB), which was the largest model fitting in memory.

\textbf{Phase 6: Optimization Attempt.} We attempted to optimize llama2:latest using the balanced priority strategy with pre- and post-optimization benchmarking enabled. The optimization process was monitored for successful completion or failure, with error messages recorded for failure analysis.

\textbf{Phase 7: Failure Case Analysis.} For failed benchmarks (qwen3:8b timeout, nomic-embed-text:latest embedding model error), we analyzed root causes by examining system logs and error messages to identify the underlying issues (memory constraints, model type incompatibility).

\subsection{Experimental Setup}

Evaluation conducted on:
\begin{itemize}
    \item \textbf{Platform}: macOS (Apple M2)
    \item \textbf{CPU}: Apple M2 (8 cores)
    \item \textbf{RAM}: 1.2 GB available
    \item \textbf{GPU}: Apple M2 (4.0 GB VRAM)
    \item \textbf{Ollama Version}: Latest
    \item \textbf{Ollama-Optimizer Version}: 1.0.0
\end{itemize}

Six models were installed:
\begin{itemize}
    \item gemma3:27b (Q4\_K\_M, 16.2 GB, 27.4B parameters)
    \item glm-5:cloud (unknown quantization, cloud-based)
    \item qwen3:8b (Q4\_K\_M, 4.9 GB, 8.2B parameters)
    \item nomic-embed-text:latest (F16, 266 MB, 137M parameters)
    \item llama2:latest (Q4\_0, 3.6 GB, 7B parameters)
    \item llama3:70b (Q4\_0, 37.2 GB, 70.6B parameters)
\end{itemize}

\subsection{Hardware Detection Results}

System profiling successfully detected:
\begin{itemize}
    \item OS: macOS
    \item CPU: Apple M2
    \item RAM: 1.2 GB available
    \item GPU: Apple M2 (4.0 GB VRAM)
    \item Flash Attention: Not supported
    \item Disk Free: 236.7 GB
\end{itemize}

Total available memory for model loading: 5.2 GB (1.2 GB RAM + 4.0 GB VRAM).

\subsection{Model Analysis Results}

Model analysis successfully parsed quantization levels for 5 models:
\begin{table}[htbp]
\centering
\caption{Installed models with quantization analysis.}
\label{tab:models}
\begin{tabular}{lccccc}
\hline
Model & Params & Quant & Size & Quality & Speed \\
\hline
gemma3:27b & 27.4B & Q4\_K\_M & 16.2 GB & 85\% & 2.3x \\
glm-5:cloud & ? & ? & 0 MB & 50\% & 1.0x \\
qwen3:8b & 8.2B & Q4\_K\_M & 4.9 GB & 85\% & 2.3x \\
nomic-embed-text:latest & 137M & F16 & 266 MB & 100\% & 1.0x \\
llama2:latest & 7B & Q4\_0 & 3.6 GB & 78\% & 2.5x \\
llama3:70b & 70.6B & Q4\_0 & 37.2 GB & 78\% & 2.5x \\
\hline
\end{tabular}
\end{table}

Parameter count parsing failed for nomic-embed-text (``137M'' string) and glm-5:cloud (cloud model with no local size).

\subsection{Optimization Recommendations}

The system correctly identified that large models exceed available memory:
\begin{itemize}
    \item gemma3:27b (16.2 GB) > 5.2 GB available: No quantization recommendation
    \item llama3:70b (37.2 GB) > 5.2 GB available: No quantization recommendation
\end{itemize}

For models fitting in memory, recommendations included:
\begin{itemize}
    \item qwen3:8b: Q4\_K\_M $\rightarrow$ IQ3\_XXS (saves 34\% memory, 1.4x speed, quality delta -0.200)
    \item nomic-embed-text:latest: F16 $\rightarrow$ IQ3\_XXS (saves 0\% memory, 3.3x speed, quality delta -0.350)
    \item llama2:latest: Q4\_0 $\rightarrow$ IQ3\_XXS (saves 23\% memory, 1.3x speed, quality delta -0.130)
\end{itemize}

\subsection{Benchmark Results}

Benchmarking of llama2:latest (7B parameters, Q4\_0, 3.6 GB) was successful:
\begin{table}[htbp]
\centering
\caption{Benchmark results for llama2:latest.}
\label{tab:benchmark}
\begin{tabular}{ccccc}
\hline
Run & Tokens/sec & TTFT (ms) & Prompt tok/s & Memory (MB) \\
\hline
1 & 21.4 & 144 & 588.9 & 0 \\
2 & 18.5 & 157 & 1132.1 & 1 \\
3 & 20.5 & 352 & 172.5 & 1 \\
4 & 20.2 & 245 & 247.8 & 6 \\
5 & 11.2 & 615 & 129.9 & 1 \\
\hline
\textbf{Mean} & \textbf{18.3} & \textbf{302} & \textbf{454.2} & \textbf{1.6} \\
\textbf{Std Dev} & \textbf{4.0} & \textbf{193} & \textbf{389.5} & \textbf{2.3} \\
\textbf{CV (\%)} & \textbf{21.9} & \textbf{63.9} & \textbf{85.8} & \textbf{143.8} \\
\hline
\end{tabular}
\end{table}

\subsection{Statistical Analysis}

We calculated standard deviation (Std Dev) and coefficient of variation (CV) for each metric to assess measurement consistency. The tokens/sec metric exhibited moderate variability (CV = 21.9\%), with values ranging from 11.2 to 21.4 tok/s. Time-to-first-token showed high variability (CV = 63.9\%) due to the substantial difference between the fastest (144 ms) and slowest (615 ms) runs, likely influenced by system load and cache effects. Prompt processing rate also showed high variability (CV = 85.8\%) due to varying prompt lengths and complexity. Memory consumption showed the highest relative variability (CV = 143.8\%) due to the small absolute values (0--6 MB), though the absolute range (6 MB) is minimal in the context of total system memory.

\subsection{Optimization Attempt}

We attempted to optimize llama2:latest using the balanced priority strategy. The pre-optimization benchmark yielded 18.5 tok/s, 1081 ms TTFT, 165.7 prompt tok/s, and 2 MB memory consumption. The system recommended Q4\_0 $\rightarrow$ IQ3\_XXS quantization, predicting 23\% memory savings (3.6 GB $\rightarrow$ 2.7 GB), 1.3x speed improvement, and -13\% quality delta.

However, the optimization failed during model creation with error: ``neither 'from' or 'files' was specified''. This indicates a bug in the Modfile generation logic where the system fails to properly specify the source model in the Ollama Modfile. The error occurred during the POST request to http://localhost:11434/api/create, suggesting the optimizer's Modfile template is incomplete or incorrectly formatted for this model.

This failure prevents actual before/after performance comparison, representing a limitation of the current implementation. The optimization planning and recommendation phases work correctly, but the model creation step requires bug fixes to generate valid Ollama Modfiles.

\subsection{Additional Benchmark Failures}

\textbf{qwen3:8b Timeout}: Benchmarking of qwen3:8b (8.2B parameters, Q4\_K\_M, 4.9 GB) failed due to timeout after 600 seconds. The model size (4.9 GB) represents 94\% of the total available memory (5.2 GB = 1.2 GB RAM + 4.0 GB VRAM), leaving only 0.3 GB for operating system overhead, KV cache, and activation storage. This insufficient margin causes excessive swapping between RAM and disk, leading to severe performance degradation and eventual timeout. The system correctly identified that this model barely fits in memory but did not prevent benchmarking attempts. This suggests the need for a safety margin threshold (e.g., 80\% of available memory) to prevent benchmarking of models that are likely to timeout.

\textbf{nomic-embed-text:latest Embedding Model Error}: Benchmarking of nomic-embed-text:latest (137M parameters, F16, 266 MB) failed with error ``does not support generate''. This model is an embedding model designed for generating text embeddings rather than text generation. The benchmark module currently assumes all models support text generation via the `/api/generate` endpoint, but embedding models require the `/api/embed` endpoint instead. This failure highlights the need for model type detection to distinguish between text generation models, embedding models, and code models, with appropriate benchmarking strategies for each type.

\subsection{Discussion}

The experimental results demonstrate:
\begin{enumerate}
    \item \textbf{Accurate hardware detection}: System profiling correctly detected CPU, RAM, GPU, and VRAM across macOS with Apple Silicon.
    \item \textbf{Successful model analysis}: Quantization levels were correctly parsed for 5 out of 6 models. Cloud-based models (glm-5:cloud) and models with non-standard parameter strings (nomic-embed-text) had parsing issues.
    \item \textbf{Correct memory constraint identification}: The system correctly identified that gemma3:27b (16.2 GB) and llama3:70b (37.2 GB) exceed available memory (5.2 GB) and appropriately warned the user.
    \item \textbf{Successful benchmarking}: llama2:latest (3.6 GB) benchmarked successfully with 18.3 tok/s throughput, demonstrating the system can measure performance on models that fit in memory.
    \item \textbf{Optimization planning works correctly}: The system correctly recommended Q4\_0 $\rightarrow$ IQ3\_XXS quantization for llama2:latest with predicted 23\% memory savings and 1.3x speed improvement.
    \item \textbf{Optimization execution bug}: The actual model creation failed due to a Modfile generation bug (``neither 'from' or 'files' was specified''), indicating the optimizer's Modfile template is incomplete or incorrectly formatted.
    \item \textbf{Timeout on borderline models}: qwen3:8b (4.9 GB) timed out during benchmarking, indicating that models barely fitting in available memory may experience performance degradation due to swapping.
    \item \textbf{Embedding model handling}: nomic-embed-text:latest failed benchmarking because embedding models do not support text generation, highlighting the need for model-type-aware benchmarking.
\end{enumerate}

The system's optimization recommendations (Q4\_K\_M $\rightarrow$ IQ3\_XXS) suggest aggressive quantization for memory-constrained environments. However, IQ3\_XXS quality retention (~50\%) may be too low for production use, suggesting the ``balanced'' priority (Q4\_K\_M at 85\% quality) is more appropriate for most users.

\section{Limitations and Future Work}

\subsection{Limitations}

\begin{enumerate}
    \item \textbf{Optimization execution bug}: The model creation step fails with ``neither 'from' or 'files' was specified'' error, indicating the Modfile generation logic is incomplete or incorrectly formatted. This prevents actual model optimization and before/after performance comparison.
    \item \textbf{Memory-constrained evaluation}: The test system had only 1.2 GB available RAM, limiting evaluation to small models (<5 GB). Large models (>10 GB) could not be benchmarked.
    \item \textbf{Model type detection}: The system does not distinguish between text generation models and embedding models, causing benchmark failures on nomic-embed-text.
    \item \textbf{Parameter count parsing}: String parsing for parameter counts (``137M'') failed in some cases.
    \item \textbf{Quality estimation}: Quality retention estimates are based on quantization level only, not actual perplexity measurements or task-specific benchmarks.
    \item \textbf{Cloud model support}: Cloud-based models (glm-5:cloud) are not fully supported as they have no local storage footprint.
\end{enumerate}

\subsection{Future Work}

\begin{enumerate}
    \item \textbf{Fix Modfile generation bug}: Debug and fix the ``neither 'from' or 'files' was specified'' error in the model creation step to enable actual optimization and before/after performance comparison.
    \item \textbf{Model type detection}: Add detection for embedding models, code models, and multimodal models with appropriate benchmarking strategies.
    \item \textbf{Perplexity measurement}: Implement actual perplexity calculation on calibration datasets to provide more accurate quality estimates.
    \item \textbf{Task-specific benchmarks}: Integrate MMLU, HellaSwag, or other standardized benchmarks for quality evaluation.
    \item \textbf{GPU acceleration}: Add GPU-accelerated quantization for faster model optimization.
    \item \textbf{Cloud optimization}: Support for optimizing models for cloud deployment with cost-aware recommendations.
    \item \textbf{Multi-model optimization}: Simultaneous optimization of multiple models considering total memory constraints.
\end{enumerate}

\section{Conclusion}

We presented Ollama-Optimizer, an automated system for hardware-aware LLM optimization that democratizes access to quantization techniques. The system successfully detects hardware capabilities, analyzes installed models, recommends optimal quantization configurations, and benchmarks performance improvements. Evaluation on an Apple M2 system demonstrated accurate hardware detection, correct memory constraint identification, and successful benchmarking of models fitting in available memory.

The system addresses the critical barrier to LLM deployment on consumer hardware by automating the complex optimization decisions traditionally requiring expert knowledge. By providing a unified, cross-platform solution with built-in quantization education, Ollama-Optimizer enables researchers, developers, and enthusiasts to deploy LLMs on their local hardware without requiring specialized expertise.

Future work will focus on improving model type detection, implementing actual quality measurements, and adding support for cloud deployment scenarios. As LLMs continue to grow in size and capability, automated optimization tools like Ollama-Optimizer will become increasingly important for democratizing access to these powerful models.

\begin{thebibliography}{99}

\bibitem{brown2020}
T. B. Brown, B. Mann, N. Ryder, M. Subbiah, J. Kaplan, P. Dhariwal, A. Girshick, et al., ``Language models are few-shot learners,'' in \textit{Adv. Neural Inf. Process. Syst.}, vol. 33, pp. 1877--1901, 2020.

\bibitem{touvron2023}
H. Touvron, L. Martin, K. Stone, P. Albert, Y. Alwahbi, A. Babaei, A. Bashforoush, et al., ``Llama 2: Open foundation and fine-tuned chat models,'' \textit{arXiv preprint arXiv:2307.09288}, Jul. 2023.

\bibitem{jiang2023}
A. Q. Jiang, A. Sablayrolles, A. Mensch, C. Bamford, D. S. Chaplot, D. de las Casas, F. Bressand, et al., ``Mistral 7B,'' \textit{arXiv preprint arXiv:2310.06825}, Oct. 2023.

\bibitem{frantar2022}
E. Frantar, S. Ashkboos, T. Hoefler, and D. Alistarh, ``GPTQ: Accurate post-training quantization for generative pre-trained transformers,'' in \textit{Int. Conf. Learn. Represent.}, arXiv:2210.17323, Oct. 2022.

\bibitem{lin2023}
J. Lin, J. Tang, H. Tang, S. Yang, X. Dang, C. Gan, and S. Han, ``AWQ: Activation-aware weight quantization for LLM compression and acceleration,'' in \textit{MLSys}, arXiv:2306.00978, Jun. 2023.

\bibitem{xiao2022}
G. Xiao, J. Lin, M. Seznec, H. Wu, J. Demouth, and S. Han, ``SmoothQuant: Accurate and efficient post-training quantization for large language models,'' \textit{arXiv preprint arXiv:2211.10438}, Nov. 2022.

\bibitem{gerganov2023}
G. Gerganov, ``llama.cpp: Port of Facebook's LLaMA model in C/C++,'' \textit{GitHub repository}, 2023. [Online]. Available: https://github.com/ggerganov/llama.cpp

\bibitem{ngrok2023}
ngrok, ``Quantization from the ground up,'' \textit{Technical Blog}, 2023. [Online]. Available: https://ngrok.com/blog/quantization

\bibitem{llamacpp2023}
ggml-org, ``llama.cpp quantization methods overview,'' \textit{GitHub Discussions}, 2023. [Online]. Available: https://github.com/ggml-org/llama.cpp/discussions/2094

\bibitem{jacob2018}
B. Jacob, S. Kligys, B. Chen, M. Zhong, Y. Ganesh, A. Bengio, and S. Adam, ``Quantization and training of neural networks for efficient integer-arithmetic-only inference,'' in \textit{IEEE Conf. Comput. Vis. Pattern Recognit.}, pp. 2708--2717, 2018.

\bibitem{liu2024}
Z. Liu, W. Wang, Y. Li, C. Zhang, T. Yang, and X. Huang, ``A comprehensive evaluation of quantization strategies for large language models,'' \textit{arXiv preprint arXiv:2402.16775}, Feb. 2024.

\end{thebibliography}

\end{document}
